\section{Simultaneous factors in genuine Jacobian neighborhoods}

\begin{lemma}[Native products and coordinate monomial products]\label{lem:sc0-native-family-monomial-product}
\lean{EnvelopeProperty.Char0.prod_span_monomial_family}
\leanok
Taking the actual native ideal on an affine open preserves finite products.
Consequently, a finite family of coordinate monomial ideals has total ideal
the coordinate monomial whose exponents are the sums of the family exponents.
\end{lemma}
\begin{proof}
The native ideal construction is multiplicative and sends the unit ideal to
the unit ideal. Evaluate its monoid homomorphism on the finite product.
Products of principal monomial ideals are principal ideals of the product;
commute the two finite products and add powers of each fixed coordinate.
\end{proof}

\begin{definition}[Actual simultaneous family Jacobian neighborhood]\label{def:sc0-snc-family-jacobian-neighbourhood}
\lean{EnvelopeProperty.Char0.SNCFamilyJacobianNeighbourhood}
\leanok
An actual adapted Jacobian neighborhood for a finite family retains the
actual affine chart, its actual \'etale coordinates and adapted boundary
components, the actual monomial formula for every family factor in those
same coordinates, and the actual principal-neighborhood Jacobian identity
for their total transform. This is literal existential geometric output.
\end{definition}

\begin{lemma}[Such simultaneous neighborhoods exist at every closed point]\label{lem:sc0-snc-family-jacobian-at-closed}
\lean{EnvelopeProperty.Char0.exists_snc_family_jacobian_neighbourhood_at_closed}
\leanok
\uses{def:sc0-snc-family-jacobian-neighbourhood,lem:sc0-native-family-monomial-product,lem:sc0-snc-simultaneous-neighbourhood,thm:sc0-native-snc-jacobian-neighbourhood}
If the actual product of the finite family is the actual SNC total transform
and the original map is isomorphic off that total ideal, every closed point
has such a simultaneous family Jacobian neighborhood in actual source
\'etale coordinates.
\end{lemma}
\begin{proof}
Use the checked actual simultaneous SNC factor chart. Its factor monomials
give the total-transform monomial by the preceding product calculation.
Apply the actual principal Jacobian-neighborhood theorem using exactly those
adapted coordinates, retaining every factor formula on the original chart.
\end{proof}

\begin{lemma}[The simultaneous neighborhoods give a genuine finite cover]\label{lem:sc0-finite-snc-family-jacobian-cover}
\lean{EnvelopeProperty.Char0.exists_finite_snc_family_jacobian_neighbourhood_cover}
\leanok
\uses{lem:sc0-snc-family-jacobian-at-closed}
On a quasi-compact actual smooth final scheme, finitely many simultaneous
family Jacobian neighborhoods cover the whole scheme.
\end{lemma}
\begin{proof}
Choose and shrink at every closed point before selecting a finite subcover.
The actual smooth finite-type scheme is Jacobson, so every point lies in
its closed specialization's chosen open neighborhood. These actual shrunk
opens therefore cover, and compactness selects finitely many of them.
\end{proof}
